% Local adaptation of Beautybook (Shilong Lu, LPPL 1.3c).
% The upstream class is provided by TeX Live and kept unmodified.
\ProvidesFile{beautybook-setup-local.tex}[2026/11/04 Discrete math book layout]

% Preserve the course fonts and the existing shared theorem counter.
\let\songti\undefined
\let\heiti\undefined
\let\kaishu\undefined
\let\fangsong\undefined
% Copyright 2018  李文威 (Wen-Wei Li).
% Permission is granted to copy, distribute and/or modify this
% document under the terms of the Creative Commons
% Attribution 4.0 International (CC BY 4.0)
% http://creativecommons.org/licenses/by/4.0/

% 目的: 字体相关设置, 呼叫相关宏包.
% 将由 AJbook.cls 引入
% 必须提供 \kaishu, \songti, \heiti, \thmheiti, \fangsong 几种字型切换命令, 在文档类中使用.
\ProvidesFile{font-setup-local.tex}[2026/10/07 Project text font files]

% 正文及等宽使用 Latin Modern，无衬线正文使用 Arial。
% no-math 保持传统 LaTeX 数学字体不被 fontspec 改写。
\setmainfont{lmroman10-regular.otf}[BoldFont=lmroman10-bold.otf,
  ItalicFont=lmroman10-italic.otf,BoldItalicFont=lmroman10-bolditalic.otf,Ligatures=TeX]
\setsansfont{Arial}[BoldFont={Arial Bold},
  ItalicFont={Arial Italic},BoldItalicFont={Arial Bold Italic},Ligatures=TeX]
\setmonofont{lmmono10-regular.otf}[ItalicFont=lmmono10-italic.otf,
  BoldFont=lmmonolt10-bold.otf,BoldItalicFont=lmmonolt10-boldoblique.otf]

% 使用用户提供的项目字体；从 notes/ 编译，源码包保留同级 font/。
% 宋体正文；黑体用于标题和粗体；楷体用于中文斜体；仿宋用于定理正文。
\setCJKmainfont[Path=../font/,
  BoldFont=FZHTK.TTF, ItalicFont=FZKTK.TTF,
  BoldItalicFont=FZHTK.TTF]{FZSSK.TTF}
\setCJKsansfont[Path=../font/,
  BoldFont=FZHTK.TTF, ItalicFont=FZHTK.TTF,
  BoldItalicFont=FZHTK.TTF]{FZHTK.TTF}
\setCJKmonofont[Path=../font/,
  BoldFont=FZHTK.TTF, ItalicFont=FZHTK.TTF,
  BoldItalicFont=FZHTK.TTF]{FZHTK.TTF}

\setCJKfamilyfont{song}[Path=../font/,
  BoldFont=FZHTK.TTF, ItalicFont=FZKTK.TTF,
  BoldItalicFont=FZHTK.TTF]{FZSSK.TTF}
\setCJKfamilyfont{kai}[Path=../font/,
  BoldFont=FZHTK.TTF, ItalicFont=FZKTK.TTF,
  BoldItalicFont=FZHTK.TTF]{FZKTK.TTF}
\setCJKfamilyfont{fangsong}[Path=../font/,
  BoldFont=FZHTK.TTF, ItalicFont=FZKTK.TTF,
  BoldItalicFont=FZHTK.TTF]{FZFSK.TTF}
\setCJKfamilyfont{hei}[Path=../font/,
  BoldFont=FZHTK.TTF, ItalicFont=FZHTK.TTF,
  BoldItalicFont=FZHTK.TTF]{FZHTK.TTF}
\setCJKfamilyfont{hei2}[Path=../font/,BoldFont=FZHTK.TTF]{FZHTK.TTF}
\setCJKfamilyfont{sectionfont}[Path=../font/,BoldFont=FZHTK.TTF]{FZHTK.TTF}
\setCJKfamilyfont{pffont}[Path=../font/,BoldFont=FZHTK.TTF]{FZHTK.TTF}
\setCJKfamilyfont{emfont}[Path=../font/,BoldFont=FZHTK.TTF]{FZHTK.TTF}

\punctstyle{kaiming}

\defaultfontfeatures{Ligatures=TeX} 
\XeTeXlinebreaklocale "zh"
\XeTeXlinebreakskip = 0pt plus 1pt minus 0.1pt

% 以下设置字体相关命令, 用于定理等环境中.
\newcommand\kaishu{\CJKfamily{kai}} % 楷体
\newcommand\songti{\CJKfamily{song}} % 宋体
\newcommand\heiti{\CJKfamily{hei}}	% 黑体
\newcommand\thmheiti{\CJKfamily{hei2}}	% 用于定理名称的黑体
\newcommand\fangsong{\CJKfamily{fangsong}} % 仿宋
\renewcommand{\em}{\bfseries\CJKfamily{emfont}} % 强调
\usepackage{amssymb,mathtools}
\usepackage[no-math,nohelv,nott]{newtxtext}
% Keep newtxtext's global OTF extension out of Arial and xeCJK font lookup.
\defaultfontfeatures{Ligatures=TeX}
\usepackage[varvw,varbb]{newtxmath}
% Default calligraphic alphabet; reserve script for distinct special notation.
\usepackage[cal=euler,scr=rsfso]{mathalpha}
% Shared math settings for the main notes and the standalone documents.
\ProvidesFile{math-setup-local.tex}[2026/11/10 Local math glyphs and conventions]
% Provide semantic double brackets.
\usepackage{stmaryrd}
% St Mary Road has no bold series; keep its bracket glyphs at their native weight.
\SetSymbolFont{stmry}{bold}{U}{stmry}{m}{n}
% AJbook loads ntheorem at the end of the preamble and defines its own QED box.
\let\openbox\relax

% Put limits beside all amsmath integral signs, also in display style.
% Override AJbook's intlimits option locally without changing the class.
\makeatletter
\let\ilimits@\nolimits
\makeatother

% Standalone documents use Computer Modern with just the varbb glyphs.
% The main notes already get these glyphs from newtxmath.
\ifdefined\newtxmathLoaded\else
\DeclareFontFamily{U}{projectvarbb}{}
\DeclareFontShape{U}{projectvarbb}{m}{n}{<-> ntxmia}{}
\DeclareFontShape{U}{projectvarbb}{b}{n}{<-> ntxbmia}{}
\DeclareSymbolFont{projectvarbb}{U}{projectvarbb}{m}{n}
\SetSymbolFont{projectvarbb}{bold}{U}{projectvarbb}{b}{n}
\ExplSyntaxOn
\exp_args:Nc \DeclareMathSymbol {project_bb_A:} {\mathord}{projectvarbb}{132}
\exp_args:Nc \DeclareMathSymbol {project_bb_B:} {\mathord}{projectvarbb}{133}
\exp_args:Nc \DeclareMathSymbol {project_bb_C:} {\mathord}{projectvarbb}{134}
\exp_args:Nc \DeclareMathSymbol {project_bb_D:} {\mathord}{projectvarbb}{135}
\exp_args:Nc \DeclareMathSymbol {project_bb_E:} {\mathord}{projectvarbb}{136}
\exp_args:Nc \DeclareMathSymbol {project_bb_F:} {\mathord}{projectvarbb}{137}
\exp_args:Nc \DeclareMathSymbol {project_bb_G:} {\mathord}{projectvarbb}{138}
\exp_args:Nc \DeclareMathSymbol {project_bb_H:} {\mathord}{projectvarbb}{139}
\exp_args:Nc \DeclareMathSymbol {project_bb_I:} {\mathord}{projectvarbb}{140}
\exp_args:Nc \DeclareMathSymbol {project_bb_J:} {\mathord}{projectvarbb}{141}
\exp_args:Nc \DeclareMathSymbol {project_bb_K:} {\mathord}{projectvarbb}{142}
\exp_args:Nc \DeclareMathSymbol {project_bb_L:} {\mathord}{projectvarbb}{143}
\exp_args:Nc \DeclareMathSymbol {project_bb_M:} {\mathord}{projectvarbb}{144}
\exp_args:Nc \DeclareMathSymbol {project_bb_N:} {\mathord}{projectvarbb}{145}
\exp_args:Nc \DeclareMathSymbol {project_bb_O:} {\mathord}{projectvarbb}{146}
\exp_args:Nc \DeclareMathSymbol {project_bb_P:} {\mathord}{projectvarbb}{147}
\exp_args:Nc \DeclareMathSymbol {project_bb_Q:} {\mathord}{projectvarbb}{148}
\exp_args:Nc \DeclareMathSymbol {project_bb_R:} {\mathord}{projectvarbb}{149}
\exp_args:Nc \DeclareMathSymbol {project_bb_S:} {\mathord}{projectvarbb}{150}
\exp_args:Nc \DeclareMathSymbol {project_bb_T:} {\mathord}{projectvarbb}{151}
\exp_args:Nc \DeclareMathSymbol {project_bb_U:} {\mathord}{projectvarbb}{152}
\exp_args:Nc \DeclareMathSymbol {project_bb_V:} {\mathord}{projectvarbb}{153}
\exp_args:Nc \DeclareMathSymbol {project_bb_W:} {\mathord}{projectvarbb}{154}
\exp_args:Nc \DeclareMathSymbol {project_bb_X:} {\mathord}{projectvarbb}{155}
\exp_args:Nc \DeclareMathSymbol {project_bb_Y:} {\mathord}{projectvarbb}{156}
\exp_args:Nc \DeclareMathSymbol {project_bb_Z:} {\mathord}{projectvarbb}{157}
\exp_args:Nc \DeclareMathSymbol {project_bb_a:} {\mathord}{projectvarbb}{158}
\exp_args:Nc \DeclareMathSymbol {project_bb_b:} {\mathord}{projectvarbb}{159}
\exp_args:Nc \DeclareMathSymbol {project_bb_c:} {\mathord}{projectvarbb}{160}
\exp_args:Nc \DeclareMathSymbol {project_bb_d:} {\mathord}{projectvarbb}{161}
\exp_args:Nc \DeclareMathSymbol {project_bb_e:} {\mathord}{projectvarbb}{162}
\exp_args:Nc \DeclareMathSymbol {project_bb_f:} {\mathord}{projectvarbb}{163}
\exp_args:Nc \DeclareMathSymbol {project_bb_g:} {\mathord}{projectvarbb}{164}
\exp_args:Nc \DeclareMathSymbol {project_bb_h:} {\mathord}{projectvarbb}{165}
\exp_args:Nc \DeclareMathSymbol {project_bb_i:} {\mathord}{projectvarbb}{166}
\exp_args:Nc \DeclareMathSymbol {project_bb_j:} {\mathord}{projectvarbb}{167}
\exp_args:Nc \DeclareMathSymbol {project_bb_k:} {\mathord}{projectvarbb}{168}
\exp_args:Nc \DeclareMathSymbol {project_bb_l:} {\mathord}{projectvarbb}{169}
\exp_args:Nc \DeclareMathSymbol {project_bb_m:} {\mathord}{projectvarbb}{170}
\exp_args:Nc \DeclareMathSymbol {project_bb_n:} {\mathord}{projectvarbb}{171}
\exp_args:Nc \DeclareMathSymbol {project_bb_o:} {\mathord}{projectvarbb}{172}
\exp_args:Nc \DeclareMathSymbol {project_bb_p:} {\mathord}{projectvarbb}{173}
\exp_args:Nc \DeclareMathSymbol {project_bb_q:} {\mathord}{projectvarbb}{174}
\exp_args:Nc \DeclareMathSymbol {project_bb_r:} {\mathord}{projectvarbb}{175}
\exp_args:Nc \DeclareMathSymbol {project_bb_s:} {\mathord}{projectvarbb}{176}
\exp_args:Nc \DeclareMathSymbol {project_bb_t:} {\mathord}{projectvarbb}{177}
\exp_args:Nc \DeclareMathSymbol {project_bb_u:} {\mathord}{projectvarbb}{178}
\exp_args:Nc \DeclareMathSymbol {project_bb_v:} {\mathord}{projectvarbb}{179}
\exp_args:Nc \DeclareMathSymbol {project_bb_w:} {\mathord}{projectvarbb}{180}
\exp_args:Nc \DeclareMathSymbol {project_bb_x:} {\mathord}{projectvarbb}{181}
\exp_args:Nc \DeclareMathSymbol {project_bb_y:} {\mathord}{projectvarbb}{182}
\exp_args:Nc \DeclareMathSymbol {project_bb_z:} {\mathord}{projectvarbb}{183}
\exp_args:Nc \DeclareMathSymbol {project_bb_0:} {\mathord}{projectvarbb}{43}
\exp_args:Nc \DeclareMathSymbol {project_bb_1:} {\mathord}{projectvarbb}{44}
\exp_args:Nc \DeclareMathSymbol {project_bb_2:} {\mathord}{projectvarbb}{45}
\exp_args:Nc \DeclareMathSymbol {project_bb_3:} {\mathord}{projectvarbb}{46}
\exp_args:Nc \DeclareMathSymbol {project_bb_4:} {\mathord}{projectvarbb}{47}
\exp_args:Nc \DeclareMathSymbol {project_bb_5:} {\mathord}{projectvarbb}{48}
\exp_args:Nc \DeclareMathSymbol {project_bb_6:} {\mathord}{projectvarbb}{49}
\exp_args:Nc \DeclareMathSymbol {project_bb_7:} {\mathord}{projectvarbb}{50}
\exp_args:Nc \DeclareMathSymbol {project_bb_8:} {\mathord}{projectvarbb}{51}
\exp_args:Nc \DeclareMathSymbol {project_bb_9:} {\mathord}{projectvarbb}{52}
\RenewDocumentCommand{\mathbb}{m}{
  \tl_map_inline:nn {#1} {\use:c {project_bb_##1:}}
}
\ExplSyntaxOff
\fi

% Use Arial only for the math sans-serif alphabet, including bold math.
\edef\projectmathsffamily{\sfdefault}
\DeclareMathAlphabet{\mathsf}{TU}{\projectmathsffamily}{m}{n}
\SetMathAlphabet{\mathsf}{bold}{TU}{\projectmathsffamily}{b}{n}

\hypersetup{unicode,bookmarksnumbered=true,hidelinks}
\setlength{\parindent}{2em}
\setcounter{secnumdepth}{2}
\renewcommand{\Re}{\operatorname{Re}}
\renewcommand{\Im}{\operatorname{Im}}
\renewcommand{\emptyset}{\ensuremath{\varnothing}}

\newtheoremstyle{coursestatement}{6pt}{6pt}{\fangsong}{}%
  {\sffamily\bfseries\thmheiti}{}{.5em}%
  {\thmname{#1}\thmnumber{ #2}\thmnote{（#3）}}
\newtheoremstyle{courseexample}{6pt}{6pt}{\normalfont}{}%
  {\sffamily\bfseries\thmheiti}{}{.5em}%
  {\thmname{#1}\thmnumber{ #2}\thmnote{（#3）}}
\theoremstyle{coursestatement}
\newtheorem{theorem}{定理}[section]
% Set the section prefix before thmtools aliases the other theorem counters.
\renewcommand{\theHtheorem}{\theHsection.\arabic{theorem}}
\newtheorem{definition}[theorem]{定义}
\newtheorem{lemma}[theorem]{引理}
\newtheorem{proposition}[theorem]{命题}
\newtheorem{corollary}[theorem]{推论}
\newtheorem{axiom}[theorem]{公理}
\newtheorem{definition-theorem}[theorem]{定义--定理}
\newtheorem{definition-proposition}[theorem]{定义--命题}
\newtheorem{hypothesis}[theorem]{假设}
\newtheorem{conjecture}[theorem]{猜想}
\theoremstyle{courseexample}
\newtheorem{example}[theorem]{例}
\newtheorem{remark}[theorem]{注记}
\newtheorem{convention}[theorem]{约定}
\newtheorem{exercise}[theorem]{练习}

\renewcommand{\qedsymbol}{\ensuremath{\square}}
\makeatletter
\renewenvironment{proof}[1][]{%
  \par\pushQED{\qed}\normalfont\topsep6pt
  \trivlist\item[\hskip\labelsep\sffamily\bfseries\heiti
    证明\ifstrempty{#1}{}{（#1）}]\ignorespaces
}{\popQED\endtrivlist\@endpefalse}
% Chapters are the top level of this book.
\patchcmd{\ttlf@chapter}{Part\ \thepart}{课程笔记}{}{%
  \PackageError{beautybook-setup-local}{Unexpected Beautybook chapter format}{}}
\makeatother

% The upstream chapter heading is a page overlay.
\geometry{left=2.15cm,right=2.15cm,top=2.4cm,bottom=2.45cm,headheight=18pt}
\setlength{\footheight}{18pt}
\titlespacing{\chapter}{0pt}{0pt}{101pt}
% A standalone heading also works before theorem boxes and lists.
\titleformat{\paragraph}[block]{\normalfont\sffamily\bfseries\heiti}{}{0pt}{#1}
\titlespacing*{\paragraph}{0pt}{1.2ex plus .3ex minus .2ex}{.6ex plus .2ex}
\tikzfading[name=fade up,bottom color=transparent!0,top color=transparent!100]
\chapimage{artwork/course-chapter-art.pdf}
\renewcommand{\sc}{\normalfont\heiti}
\patchcmd{\mysectionformat}{Sec\ \thesection}{\S\ \thesection}{}{%
  \PackageError{beautybook-setup-local}{Unexpected Beautybook section format}{}}

% Lighten the full-height stripe and remove its checkerboard.
% Keep the page number, chapter badge, and chapter opener at the original contrast.
\colorlet{coursepageband}{outermarginfgcolor!6}
\ModifyLayer[contents={\color{coursepageband}\rule{\layerwidth}{\layerheight}}]{scth.outermargin.bg.even}
\ModifyLayer[contents={\color{coursepageband}\rule{\layerwidth}{\layerheight}}]{scth.outermargin.bg.odd}
\RemoveLayersFromPageStyle{scrheadings}{scth.outermargin.fg.even,scth.outermargin.fg.odd}
\RemoveLayersFromPageStyle{plain.scrheadings}{scth.outermargin.fg.even,scth.outermargin.fg.odd}

\newtcolorbox{wenxintishi}{enhanced,breakable,sharp corners,
  colback=outermarginbgcolor!15,colframe=outermarginfgcolor!70,boxrule=.35pt,
  colbacktitle=outermarginbgcolor!15,coltitle=black,
  fonttitle=\sffamily\bfseries\heiti,titlerule=0pt,
  title={阅读提示},before skip=9pt,after skip=9pt}

\tcbset{coursebeauty/.style={enhanced,breakable,sharp corners,
    boxrule=0pt,frame hidden,borderline west={1.15pt}{0pt}{outermarginfgcolor},
    left=1mm,right=1mm,top=.5mm,bottom=.5mm,boxsep=1mm,
    before skip=6pt,after skip=6pt}}
\tcolorboxenvironment{definition}{coursebeauty,colback=outermarginbgcolor!24}
\tcolorboxenvironment{theorem}{coursebeauty,colback=outermarginbgcolor!40}
\tcolorboxenvironment{lemma}{coursebeauty,colback=outermarginbgcolor!30}
\tcolorboxenvironment{proposition}{coursebeauty,colback=outermarginbgcolor!24}
\tcolorboxenvironment{corollary}{coursebeauty,colback=outermarginbgcolor!24}

\newcommand{\makecoursecover}[1]{%
  \begin{titlepage}
  \thispagestyle{empty}
  \begin{tikzpicture}[remember picture,overlay]
  \fill[outermarginbgcolor!12] (current page.south west) rectangle (current page.north east);
  \fill[outermarginbgcolor!42] (current page.north west)
    rectangle ([yshift=-100mm]current page.north east);
  \fill[outermarginfgcolor] (current page.north west)
    rectangle ([yshift=-3mm]current page.north east);
  \node[anchor=north west,inner sep=0,text=outermarginfgcolor,
    font=\sffamily\bfseries\fontsize{11}{14}\selectfont]
    at ([xshift=22mm,yshift=-23mm]current page.north west) {课程笔记\quad /\quad COURSE NOTES};
  \node[anchor=north east,inner sep=0,text=outermarginfgcolor,
    font=\sffamily\fontsize{23}{26}\selectfont]
    at ([xshift=-22mm,yshift=-20mm]current page.north east) {2026};
  \node[anchor=north west,inner sep=0,
    font=\sffamily\bfseries\fontsize{34}{44}\selectfont]
    at ([xshift=22mm,yshift=-51mm]current page.north west) {离散数学与结构};
  \node[anchor=north west,inner sep=0,text=outermarginfgcolor,
    font=\sffamily\fontsize{12}{16}\selectfont]
    at ([xshift=22mm,yshift=-73mm]current page.north west) {DISCRETE MATHEMATICS AND STRUCTURES};
  \draw[outermarginfgcolor,line width=1pt]
    ([xshift=22mm,yshift=-88mm]current page.north west) -- ++(25mm,0);

  % A geometric network echoes the chapter artwork at cover scale.
  \begin{scope}[shift={([xshift=105mm,yshift=-143mm]current page.north west)},
    x=1mm,y=1mm]
    \draw[outermarginfgcolor!10,line width=.5pt] (0,0) circle (43);
    \draw[outermarginfgcolor!10,line width=.5pt] (-52,0)--(52,0) (0,-47)--(0,47);
    \coordinate (a) at (-68,-6);
    \coordinate (b) at (-42,31);
    \coordinate (c) at (-32,-32);
    \coordinate (d) at (0,5);
    \coordinate (e) at (28,36);
    \coordinate (f) at (40,-25);
    \coordinate (g) at (69,13);
    \fill[outermarginfgcolor!5] (a)--(b)--(d)--(c)--cycle;
    \fill[outermarginfgcolor!9] (d)--(e)--(g)--(f)--cycle;
    \draw[outermarginfgcolor!55,line width=.8pt]
      (a)--(b)--(c)--cycle (b)--(d)--(c)
      (d)--(e)--(f)--cycle (e)--(g)--(f);
    \draw[outermarginfgcolor!25,line width=.6pt] (b)--(e) (c)--(f);
    \foreach \p in {a,b,c,d,e,f,g}
      \fill[outermarginfgcolor!80] (\p) circle (1.1);
    \foreach \p in {a,c,e,g}
      \draw[outermarginfgcolor!25,line width=.5pt] (\p) circle (2.6);
  \end{scope}

  % First isomorphism theorem: quotienting identifies precisely the kernel.
  \begin{scope}[shift={([xshift=35mm,yshift=-247mm]current page.north west)},
    x=1mm,y=1mm,>=Stealth,text=outermarginfgcolor!85,
    every node/.style={font=\large}]
    \node (coverG) at (0,28) {$G$};
    \node (coverIm) at (48,28) {$\operatorname{im}\varphi$};
    \node (coverQ) at (0,0) {$G/\ker\varphi$};
    \draw[->,outermarginfgcolor!65,line width=.7pt]
      (coverG)--node[above,font=\small] {$\varphi$} (coverIm);
    \draw[->,outermarginfgcolor!65,line width=.7pt]
      (coverG)--node[left,font=\small] {$\pi$} (coverQ);
    \draw[->,outermarginfgcolor!65,line width=.7pt]
      (coverQ)--node[below right,font=\small] {$\overline\varphi\;\cong$} (coverIm);
  \end{scope}

  % A genuine Poisson(3) probability mass function, rather than arbitrary bars.
  \begin{scope}[shift={([xshift=125mm,yshift=-247mm]current page.north west)},
    x=1mm,y=1mm,>=Stealth,text=outermarginfgcolor!85]
    \node[anchor=west,font=\large] at (0,39) {$X\sim\operatorname{Poisson}(3)$};
    \foreach \k/\p in {0/.049787,1/.149361,2/.224042,3/.224042,4/.168031,5/.100819,6/.050409,7/.021604,8/.008102} {
      \fill[outermarginfgcolor!18] (\k*6,0) rectangle ++(3.6,\p*115);
      \draw[outermarginfgcolor!60,line width=.45pt] (\k*6,0)--++(0,\p*115)--++(3.6,0)--++(0,-\p*115);
      \node[below,font=\scriptsize] at (\k*6+1.8,-1) {$\k$};
    }
    \draw[->,outermarginfgcolor!45,line width=.6pt] (-2,0)--(57,0) node[right,font=\small] {$k$};
    \draw[->,outermarginfgcolor!45,line width=.6pt] (-2,0)--(-2,31);
    \node[anchor=east,font=\scriptsize] at (-3,29) {$P_X(k)$};
  \end{scope}
  \draw[outermarginfgcolor!35,line width=.5pt]
    ([xshift=22mm,yshift=-266mm]current page.north west) --
    ([xshift=-22mm,yshift=-266mm]current page.north east);
  \node[anchor=north west,inner sep=0,font=\small]
    at ([xshift=22mm,yshift=-273mm]current page.north west) {#1};
  \node[anchor=north east,inner sep=0,text=outermarginfgcolor,font=\sffamily\small]
    at ([xshift=-22mm,yshift=-273mm]current page.north east) {定义\;·\;证明\;·\;例题};
  \end{tikzpicture}
  \null
  \end{titlepage}%
}
